\documentclass{article}

\usepackage{todonotes}
\usepackage{geometry}
\usepackage{amsmath}
\usepackage{amssymb}
\usepackage{amsthm}
\usepackage{bm}
\usepackage{graphicx}
\graphicspath{{./figures/}}
\usepackage{color}
\usepackage{tabu}
\usepackage{booktabs}
\usepackage{caption}
\usepackage{authblk}
\usepackage{cite}
% \usepackage[numbers, sort&compress]{natbib}
\usepackage{subfigure}
\usepackage{enumerate}
\usepackage{tabu}
\usepackage{multirow}
\usepackage{diagbox}
\usepackage[ruled, linesnumbered]{algorithm2e}
\newtheorem{remark}{Remark}
\newtheorem{defi}{Definition}

\newcommand{\eps}{\varepsilon}
\newcommand{\R}{\mathbb{R}}
\newcommand{\I}{\mathbb{I}}
\newcommand{\sS}{\mathbb{S}}
\newcommand{\bn}{\bm{n}}
\newcommand{\bx}{\bm{x}}
\newcommand{\by}{\bm{y}}
\newcommand{\bz}{\bm{z}}
\newcommand{\bw}{\bm{w}}
\newcommand{\bv}{\bm{v}}
\newcommand{\bmu}{\bm{\mu}}
\newcommand{\bsigma}{\bm{\sigma}}
\newcommand{\bA}{\bm{A}}
\newcommand{\LO}{\mathcal{L}}
\newcommand*{\diff}{\mathop{}\!\mathrm{d}}
\newcommand*{\Diff}[1]{\mathop{}\!\mathrm{d^#1}}
\newcommand{\avg}[1]{\left \langle #1 \right \rangle}
\newcommand{\lb}[1]{\left ( #1 \right )}

\newcommand{\KW}[1]{\textcolor{purple}{{KW: #1}}}

\begin{document}
\title{Light-inspired Neural Operators for Solving Parametric Partial Differential Equations}

\author[]{Keke Wu}

\date{\today}

\maketitle

\section{Method}

\subsection{Operator learning framework}

Let $\Omega \subset \mathbb{R}^2$ be a bounded spatial domain, and consider a parametric family of partial differential equations characterized by a nonlinear solution operator
\begin{equation}
\mathcal{G} : \mathcal{X} \to \mathcal{Y}, 
\qquad
a \mapsto u = \mathcal{G}(a),
\end{equation}
where $a : \Omega \to \mathbb{R}^{d_{\mathrm{in}}}$ denotes the input field and $u : \Omega \to \mathbb{R}^{d_{\mathrm{out}}}$ denotes the corresponding solution. The objective is to construct a parametric approximation $\mathcal{G}_\theta$ such that
\begin{equation}
\mathcal{G}_\theta(a)(x) \approx \mathcal{G}(a)(x),
\qquad x \in \Omega.
\end{equation}

Following the neural operator paradigm, we introduce a latent representation
\begin{equation}
h : \Omega \to \mathbb{R}^M,
\end{equation}
and define the proposed operator in the form
\begin{equation}
\mathcal{G}_\theta
=
\Pi_\theta
\circ
\left(
\mathcal{B}_\theta^{(L)}
\circ \cdots \circ
\mathcal{B}_\theta^{(1)}
\right)
\circ
\mathcal{P}_\theta,
\label{eq:lno-global}
\end{equation}
where $\mathcal{P}_\theta$ and $\Pi_\theta$ denote the lifting and projection operators, respectively, and each $\mathcal{B}_\theta^{(\ell)}$ is an evolution operator acting on the latent field. The architecture is implemented on a discrete grid, but its construction is intended to approximate an underlying continuous operator.

\subsection{Discrete representation}

Let $\{x_i\}_{i=1}^N \subset \Omega$ denote a uniform discretization of the domain with $N = HW$, where $H$ and $W$ are the grid resolutions along the two spatial directions. At the discrete level, the latent field is represented by
\begin{equation}
h_i^{(\ell)} \approx h^{(\ell)}(x_i) \in \mathbb{R}^M,
\qquad i = 1,\dots,N,
\end{equation}
or equivalently by a tensor of shape $H \times W \times M$ for each sample in the batch.

Each block updates the latent state through a residual evolution law
\begin{equation}
h^{(\ell+1)}
=
h^{(\ell)}
+
\mathcal{E}_\theta(h^{(\ell)}),
\qquad i=1,\dots,N,
\label{eq:block-update}
\end{equation}
where $\mathcal{E}_\theta$ denotes the increment operator. The central modeling principle of the proposed method is to decompose $\mathcal{E}_\theta$ into a superposition of three mechanisms inspired by light transport, namely reflection, refraction, and scattering. More precisely, we write
\begin{equation}
\mathcal{E}_\theta
=
\alpha_r \mathcal{R}_\theta
+
\alpha_t \mathcal{T}_\theta
+
\alpha_s \mathcal{S}_\theta,
\label{eq:evolution-decomposition}
\end{equation}
where $\mathcal{R}_\theta$ and $\mathcal{T}_\theta$ act pointwise in feature space, whereas $\mathcal{S}_\theta$ introduces nonlocal coupling across the physical domain. The coefficients $\alpha_r$, $\alpha_t$, and $\alpha_s$ are adaptive weights satisfying
\begin{equation}
\alpha_r + \alpha_t + \alpha_s = 1,
\qquad
\alpha_r,\alpha_t,\alpha_s \ge 0.
\end{equation}

\subsection{Feature-space light transformations}

The first two components of \eqref{eq:evolution-decomposition} are local transformations with respect to the spatial variable. At each grid point $x_i$, they act on the latent feature vector $h_i \in \mathbb{R}^M$ through an adaptive direction determined by the current state. Specifically, we define
\begin{equation}
v_i = \sigma (w h_i + b),
\qquad
\widehat{v}_i = \frac{v_i}{\|v_i\| + \varepsilon},
\label{eq:adaptive-direction}
\end{equation}
where $\varepsilon > 0$ is a small regularization parameter introduced to avoid numerical singularity.

The reflection operator is defined by
\begin{equation}
\mathcal{R}_\theta(h_i)
=
h_i - 2 \langle h_i, \widehat{v}_i \rangle \widehat{v}_i.
\label{eq:reflection}
\end{equation}
This is a Householder-type transformation in the latent feature space. It performs a direction-dependent redistribution of feature components while preserving the overall geometric scale of the representation up to the regularization introduced in \eqref{eq:adaptive-direction}. Since \eqref{eq:reflection} acts independently at each spatial location, it does not introduce any spatial interaction.

The refraction operator is constructed from the decomposition of $h_i$ into components orthogonal and parallel to the adaptive direction $\widehat{v}_i$. Writing
\begin{equation}
h_i^{\perp}
=
\langle h_i, \widehat{v}_i \rangle \widehat{v}_i,
\qquad
h_i^{\parallel}
=
h_i - h_i^{\perp},
\end{equation}
we define
\begin{equation}
\mathcal{T}_\theta(h_i)
=
h_i^{\parallel}
+
\eta \, h_i^{\perp},
\label{eq:refraction}
\end{equation}
where $\eta$ is a learnable scalar parameter. In practice, $\eta$ is constrained to remain in a neighborhood of unity, so that \eqref{eq:refraction} realizes a controlled anisotropic deformation of the latent representation. This mechanism modifies the relative contribution of different feature directions without altering the locality structure of the operator.

The use of \eqref{eq:reflection} and \eqref{eq:refraction} leads to an explicit separation between spatial transport and latent feature transformation. In particular, the model allows local mixing and directional reweighting to be performed directly in feature space before nonlocal propagation is applied.

\subsection{Scattering operator as a nonlinear kernel map}

The third component in \eqref{eq:evolution-decomposition} accounts for spatial propagation. Unlike the feature-space transformations introduced above, it couples different spatial locations through an adaptive kernel. At the discrete level, the scattering operator is written as
\begin{equation}
(\mathcal{S}_\theta h)_i
=
\mu
\left(
\sum_{j=1}^N K_{ij} h_j - h_i
\right),
\qquad i=1,\dots,N,
\label{eq:scattering-discrete}
\end{equation}
where $\mu > 0$ is a learnable scaling factor and $K_{ij}$ is a normalized interaction kernel.

To define $K_{ij}$, we first introduce feature-dependent query and key embeddings
\begin{equation}
q_i = W_q h_i,
\qquad
k_j = W_k h_j,
\end{equation}
where $W_q, W_k \in \mathbb{R}^{d \times M}$ are learnable matrices. The interaction weight between locations $x_i$ and $x_j$ is then defined as
\begin{equation}
K_{ij}
=
\frac{
\exp\!\left(
\frac{q_i^\top k_j}{\sqrt{d}} + r(x_i-x_j)
\right)
}{
\sum_{j=1}^N
\exp\!\left(
\frac{q_i^\top k_{j}}{\sqrt{d}} + r(x_i-x_{j})
\right)
},
\label{eq:kernel-discrete}
\end{equation}
where the positional bias is given by
\begin{equation}
r(x_i-x_j)
=
-\tau \|x_i-x_j\|^2
\label{eq:pos-bias}
\end{equation}
with $\tau > 0$ being a learnable parameter. The Gaussian-type bias \eqref{eq:pos-bias} introduces an explicit locality prior, while the feature similarity term in \eqref{eq:kernel-discrete} allows the interaction pattern to adapt to the current latent state. Since \eqref{eq:kernel-discrete} is normalized row-wise, the resulting kernel is stochastic:
\begin{equation}
\sum_{j=1}^N K_{ij} = 1,
\qquad i=1,\dots,N.
\end{equation}

The operator \eqref{eq:scattering-discrete} may be viewed as a nonlinear diffusion-type or relaxation-type update. Indeed, it takes the form of a kernel average minus the identity contribution, which is reminiscent of collision-scattering structures in transport models and of graph Laplacian-type smoothing, while remaining data-dependent through the adaptive kernel.

From the operator-learning viewpoint, \eqref{eq:scattering-discrete} admits the following continuum interpretation:
\begin{equation}
(\mathcal{S}_\theta h)(x)
=
\mu
\left(
\int_\Omega
K_\theta(x,y;h)\, h(y)\, dy
-
h(x)
\right),
\label{eq:scattering-continuous}
\end{equation}
where $K_\theta(x,y;h)$ denotes a nonlinear kernel depending on both the spatial variables and the latent field itself. This representation makes clear that the proposed architecture belongs to the class of learned nonlinear integral operators, while preserving a structured decomposition between local feature modulation and nonlocal spatial interaction.

\subsection{Adaptive coupling of light components}

The coefficients in \eqref{eq:evolution-decomposition} are not fixed a priori. Instead, they are dynamically determined from a global summary of the latent state. Let
\begin{equation}
\overline{h}
=
\frac{1}{N}
\sum_{i=1}^N h_i
\in \mathbb{R}^M
\label{eq:global-pooling}
\end{equation}
denote the spatial average of the latent representation. We define
\begin{equation}
(\alpha_r,\alpha_t,\alpha_s)
=
\mathrm{softmax}\!\bigl(g_\theta(\overline{h})\bigr),
\label{eq:adaptive-weights}
\end{equation}
where $g_\theta : \mathbb{R}^M \to \mathbb{R}^3$ is a learnable map. This mechanism allows the model to adaptively balance the contributions of reflection, refraction, and scattering according to the input-dependent global context. In this way, the architecture is able to modulate the relative emphasis on local feature transformation and nonlocal propagation across different problem instances.

\subsection{Full evolution block}

Combining the previous components, the $\ell$-th latent update is written as
\begin{equation}
\widetilde{h}^{(\ell)}
=
\alpha_r \mathcal{R}_\theta(h^{(\ell)})
+
\alpha_t \mathcal{T}_\theta(h^{(\ell)})
+
\alpha_s \mathcal{S}_\theta(h^{(\ell)}),
\label{eq:mixture}
\end{equation}
followed by a residual projection and a pointwise nonlinear correction. More precisely, the full block takes the form
\begin{equation}
h^{(\ell+1)}
=
h^{(\ell)}
+
W_o \widetilde{h}^{(\ell)}
+
\mathcal{F}_\theta(h^{(\ell)}),
\label{eq:full-block}
\end{equation}
where $W_o \in \mathbb{R}^{M \times M}$ is a learnable linear projection and $\mathcal{F}_\theta$ is a pointwise multilayer perceptron acting on the feature dimension. In practice, normalization layers are applied before the light transformations and before the nonlinear correction, which stabilizes training and improves the conditioning of the latent dynamics.

Equation \eqref{eq:full-block} may be interpreted as a residual discretization of an abstract evolution equation in latent function space. Under this viewpoint, the feature-space transformations \eqref{eq:reflection}--\eqref{eq:refraction} serve as local state-modulation mechanisms, while the scattering term \eqref{eq:scattering-discrete} supplies nonlocal propagation. The network depth therefore plays the role of an artificial evolution horizon over which information is progressively transformed and transported.

\subsection{Lifting and projection operators}

The lifting operator maps the input field into the latent feature space through a pointwise linear transformation,
\begin{equation}
h^{(0)}(x_i) = \mathcal{P}_\theta(u)(x_i) = W_{\mathrm{in}} u(x_i) + b_{\mathrm{in}},
\label{eq:lifting}
\end{equation}
where $W_{\mathrm{in}} \in \mathbb{R}^{M \times d_{\mathrm{in}}}$ and $b_{\mathrm{in}} \in \mathbb{R}^M$ are learnable parameters. After $L$ evolution blocks, the final output is obtained by a pointwise projection
\begin{equation}
y_\theta(x_i)
=
\Pi_\theta(h^{(L)})(x_i)
=
W_{\mathrm{out}} h^{(L)}(x_i) + b_{\mathrm{out}},
\label{eq:projection}
\end{equation}
with $W_{\mathrm{out}} \in \mathbb{R}^{d_{\mathrm{out}} \times M}$ and $b_{\mathrm{out}} \in \mathbb{R}^{d_{\mathrm{out}}}$.

Combining \eqref{eq:lno-global}, \eqref{eq:block-update}, and \eqref{eq:full-block}, the overall model defines a neural operator that alternates between local feature transformation and nonlocal kernel propagation, thereby providing a structured approximation class for nonlinear solution operators.

\subsection{Architectural interpretation}

The proposed construction differs from standard neural operators in that the latent evolution is explicitly decomposed into two distinct classes of mechanisms. The reflection and refraction components act only along the latent feature dimension and therefore perform local state transformation without modifying the spatial coupling pattern. In contrast, the scattering component acts across the physical domain through a learned kernel, thereby capturing long-range spatial dependence. This separation is consistent with the intuition that many PDE solution operators involve both local constitutive transformations and nonlocal transport effects. The resulting architecture combines these two aspects within a single end-to-end trainable framework, while retaining a mathematically transparent operator form.



% \bibliographystyle{plain}
% \bibliographystyle{unsrt} % cite in the order of appearance
% \bibliography{references}
\end{document}